\section{Discussion and Threats to Validity}
\label{sec:discussion}

\paragraph*{Implications for the state of the art.}
Most configurable-system sampling work pursues \emph{cost-effectiveness}, a useful
accuracy at an affordable measurement effort, rather than the lowest possible
error. Cost-efficient, distance-based, $t$-wise, and active-learning strategies are
justified by spending less for comparable accuracy, yet that effort is almost
always reported as a sample count, and RQ1 shows the count is not the budget that
is spent. A cost-effectiveness claim stated in counts is therefore unverified, and
it can be wrong: in RQ2, at a fixed measured budget the set of statistically best
samplers changes in $38$--$55\%$ of the replicated comparisons and on every Linux
release, and the sampler the literature credits is often not the cost-effective
one. Solver-based $t$-wise sampling, recommended on x264, selects expensive
configurations and loses once time is charged, while \texttt{random}, treated as a
weak baseline, is the most over-credited by count. This is not a footnote: part of
the state of the art recommends samplers that are not as cost-effective as claimed,
so a cost or efficiency claim should henceforth be established against measured
cost rather than assumed from $n$.

\paragraph*{Navigating the cost--accuracy trade-off.}
Once cost is measured, the choice is a front, not a point. Strategy--budget
combinations form a Pareto set in the accuracy--cost plane, and the RQ2
measured-cost front differs from the count-based one in $73$--$100\%$ of scenarios. The practical
decision is to pick the operating point that reaches the accuracy a study needs at
the least measured cost, rather than fixing $n$ and reporting whatever accuracy
follows. A few choices strictly dominate, being cheaper \emph{and} more accurate,
and are unambiguous wins. The rest are genuine trade-offs, and the value of
measuring cost is that the practitioner can see and navigate them instead of
having them collapsed behind a common count.

\paragraph*{Guidance for researchers and practitioners.}
The reporting change is lightweight. State the cost dimension, report both sample
size $n$ and measured cost $\TrainCost$, and show whether the two budgets agree
through the RQ1 diagnostics, cost spread and per-strategy $\Ratio$, since a global
average hides a consistently cheap or expensive strategy. If costs are homogeneous
and ratios stay near $1$, count is a defensible proxy, the validated regime of
RQ3. Otherwise, recompute the
accuracy--cost ranking and Pareto front under measured cost and label which
conclusions are count-based and which are measured-cost based. The opportunity is
concrete and now quantified: count-based sampling leaves measurable
cost-effectiveness on the table, and \texttt{pilotCost} already recovers part of
it, reaching better accuracy at equal measured budget on several systems. It can
fail where cheap regions are uninformative, but the gap we measure, largest under
heavy-tailed costs, is a clear motivation and a sizable target for the cost-aware
samplers that future work can design.
